\VerdierPhysicalPageAnchor{0078}%
\VerdierSemanticAnchor{ch1-2-5-4}
\noindent\textbf{2.5.4.}\quad
Les opérations $\star$ et ``somme des homotopies'' possèdent les propriétés
suivantes :

\VerdierSemanticAnchor{ch1-2-5-4-eq1}
\begin{equation}
  \begin{aligned}
    u\star(b+a) &= u\star b+u\star a,\\
    (b+a)\star v &= b\star v+a\star v,\\
    uu'\star a &= u\star(u'\star a),\\
    a\star vv' &= (a\star v)\star v'.
  \end{aligned}
  \tag{2.5.4.1}
\end{equation}

\VerdierSemanticAnchor{ch1-2-5-5}
\noindent\textbf{2.5.5.}\quad
Deux morphismes $f,g\colon X\rightrightarrows Y$ sont dits
\emph{homotopes} s'il existe une homotopie qui les relie.

\VerdierSemanticAnchor{ch1-2-5-5-difference-criterion}
Deux morphismes $f,g$ sont homotopes si le morphisme $f-g$ est homotope à
zéro.

\VerdierSemanticAnchor{ch1-2-5-5-shift-stability}
Soit $f\colon X\longrightarrow Y$ un morphisme homotope à zéro. Pour tout
$\sigma\in\mathbb Z[n]$, le morphisme $T(\sigma)f$ est homotope à zéro.

\VerdierSemanticAnchor{ch1-2-5-5-null-homotopic-ideal}
L'ensemble des morphismes homotopes à zéro forme un idéal de
$\operatorname{comp}^{n}(C)$, \emph{i.e.} le composé dans les deux sens d'un
morphisme homotope à zéro avec un morphisme quelconque est homotope à zéro et
la somme directe de deux morphismes homotopes à zéro est homotope à zéro.

\VerdierSemanticAnchor{ch1-2-5-6}
\noindent\textbf{2.5.6.}\quad
Soit $\varphi\colon[n]\longrightarrow[m]$ une application.

\VerdierSemanticAnchor{ch1-2-5-6-integration-preserves-homotopy}
Les foncteurs $\int_\varphi^\Sigma$ et $\int_\varphi^\Pi$ transforment les
couples de morphismes homotopes de $\operatorname{comp}^{n}(C)$ en couples de
morphismes homotopes de $\operatorname{comp}^{m}(C)$.

\VerdierSemanticAnchor{ch1-2-5-6-multiadditive-setup}
Soient de même
$F\colon\prod_{i\in[n]}C_i\longrightarrow C$ un foncteur multiadditif,

\VerdierSemanticAnchor{ch1-2-5-6-fn-functor}
\begin{equation*}
  F^{[n]}\colon
  \prod_{i\in[n]}\operatorname{comp}(C_i)
  \longrightarrow \operatorname{comp}^{n}(C)
\end{equation*}
le foncteur décrit en (2.4.2.4), $X_i$ des objets de
$\operatorname{comp}(C_i)$ $(i\ne j)$ et
$f,g\colon X\rightrightarrows Y$ deux morphismes homotopes de
$\operatorname{comp}(C_j)$.

\VerdierSemanticAnchor{ch1-2-5-6-fn-preserves-homotopy}
Les morphismes de $\operatorname{comp}^{n}(C)$ :
\begin{equation*}
  F^{[n]}(X_1,\ldots,f,\ldots,X_n)
  \quad\text{et}\quad
  F^{[n]}(X_1,\ldots,g,\ldots,X_n)
\end{equation*}
sont homotopes.

\VerdierSemanticAnchor{ch1-2-5-6-totalizations-preserve-homotopy}
Par suite, les morphismes de $\operatorname{comp}(C)$ :
\begin{equation*}
  \operatorname{comp}_{\Pi}F(X_1,\ldots,f,\ldots,X_n)
  \quad\text{et}\quad
  \operatorname{comp}_{\Pi}F(X_1,\ldots,g,\ldots,X_n)
\end{equation*}
sont homotopes (même résultat pour $\operatorname{comp}_{\Sigma}F$).

\VerdierSemanticAnchor{ch1-2-5-6-reference}
Nous laissons au lecteur le soin de vérifier toutes ces assertions ou de se
reporter à [1].

\VerdierSemanticAnchor{ch1-2-5-7}
\noindent\textbf{2.5.7.}\quad
Désignons par $K(C)$ la catégorie suivante :

\VerdierSemanticAnchor{ch1-2-5-7-objects}
\noindent\hspace*{2em}--\ Les objets de $K(C)$ sont les objets de $\operatorname{comp}(C)$.
